%%%%%%%%%%%%%%%%%%%%%%%% NIPS package imports. %%%%%%%
\usepackage[utf8]{inputenc} % allow utf-8 input
\usepackage[T1]{fontenc}    % use 8-bit T1 fonts
% \usepackage[colorlinks,urlcolor=cyan,citecolor=blue,linkcolor=magenta]{hyperref}
\usepackage{url}            % simple URL typesetting
\usepackage{booktabs}       % professional-quality tables
\usepackage{amsfonts}       % blackboard math symbols
\usepackage{microtype}      % microtypography

\usepackage{nicefrac} 
\usepackage{subcaption}
\captionsetup[subfigure]{font=footnotesize,
              labelfont=up,
              %%singlelinecheck=false,
              %%format=hang,
              justification=centering % <-- new
             }

%%%%%%%%%%%%%%%%%%%%%%%% Custom imports. %%%%%%%

\usepackage{amsmath}
\usepackage{amsthm}
\usepackage{amssymb}
\usepackage{thmtools,thm-restate}
\usepackage{thm-restate}
\usepackage{paralist}
\usepackage{comment}
\usepackage{dsfont}
%\usepackage{natbib}

\usepackage{bbm}
\usepackage{float}
\usepackage{balance}
\usepackage{stmaryrd}

\usepackage{multicol}
\usepackage[export]{adjustbox}
\usepackage{multirow}
\usepackage{tabularx}
\usepackage{mathtools}
\usepackage{mdframed}

\usepackage{enumitem}
\setitemize{noitemsep,topsep=0pt,parsep=0pt,partopsep=0pt}




%%%%%%%%%%%%%%%%%%%%%%%%%%%%%%%% captions %%%%%%%%%%%%%%%%%%%%%%%%%%%%%%%%%%


%%%%%%%%%%%%%%%%%%%%Francy algorithm coloring  %%%%%%%%%%%%%%%%%%%%%%%%%%%

\usepackage[ruled, linesnumbered]{algorithm2e} %
\usepackage{xcolor} % 
\usepackage{tikz}
\usetikzlibrary{fit,calc}
%define a marking command
\newcommand*{\tikzmk}[1]{\tikz[remember picture,overlay,] \node (#1) {};\ignorespaces}
%define a boxing command, argument = colour of box
\newcommand{\boxit}[1]{\tikz[remember picture,overlay]{\node[yshift=3pt,fill=#1,opacity=.25,fit={(A)($(B)+(.90\linewidth,.8\baselineskip)$)}] {};}\ignorespaces}

\newcommand{\boxitshort}[1]{\tikz[remember picture,overlay]{\node[yshift=3pt,fill=#1,opacity=.25,fit={(C)($(D)+(.87\linewidth,.8\baselineskip)$)}] {};}\ignorespaces}

%define some colours according to algorithm parts (or any other method you like)
\colorlet{pink}{red!40}
\colorlet{lightblue}{blue!30}
\colorlet{lightgreen}{green!30}

%Change the separation between the columns in the matrix.
\setlength\arraycolsep{2pt}

%This is so that the algorithm does not contain a semicolon in algorithm2e
\DontPrintSemicolon

\usepackage[margin=1in]{geometry}

\usepackage{natbib}
\setcitestyle{authoryear,round,citesep={;},aysep={,},yysep={;}}
\renewcommand{\cite}[1]{\citep{#1}}
\newcommand{\redTodo}[1]{\textcolor{red}{[#1]}}
\newcommand{\CS}[1]{\textcolor{pink}{#1}}
%\newcommand{\CS}[1]{#1}

\usepackage[vvarbb]{newtxmath}


%%%%%%%%%%%%%%%%%%%%%%% constants %%%%%%%%%%%%%%%%%%%%%%%%%%%%%%%%%%%%%%%%%%%

\def\fwt{{S}}
\def\db{{D}}





%%%%%%%%%%%%%%%%%%%%%%%%%%% Algos %%%%%%%%%%%%%%%%%%%%%%%%%%%%%%%%%%%%%%%%
\def\agd{{\textnormal{\texttt{AGD}}}}
\def\abm{{\textnormal{\texttt{ABM}}}}
\def\bbabm{{\textnormal{\texttt{BB-ABM}}}}
\def\afw{{\textnormal{\texttt{AFW}}}}
\def\agdavi{{\textnormal{\texttt{AGDAVI}}}}
\def\avi{{\textnormal{\texttt{AVI}}}}
\def\bpfw{{\textnormal{\texttt{BPFW}}}}
\def\cgavi{{\textnormal{\texttt{CGAVI}}}}
\def\cg{{\textnormal{\texttt{CG}}}}
\def\fw{{\textnormal{\texttt{FW}}}}
\def\ihb{{\textnormal{\texttt{IHB}}}}
\def\ihb{{\textnormal{\texttt{IHB}}}}
\def\oavi{{\textnormal{\texttt{OAVI}}}}
\def\vca{{\textnormal{\texttt{VCA}}}}
\def\pfw{{\textnormal{\texttt{PFW}}}}
\def\svd{{\textnormal{\texttt{SVD}}}}
\def\svm{{\textnormal{\texttt{SVM}}}}
\def\cvxoracle{{\textnormal{\texttt{ORACLE}}}}
\def\wihb{{\textnormal{\texttt{WIHB}}}}



\def\pcg{{\textnormal{\texttt{PCG}}}}
\def\bpcg{{\textnormal{\texttt{BPCG}}}}




%%%%%%%%%%%%%%%%%%%%%%%% Math Operators %%%%%%%%%%%%%%%%%%%%%%%%%%%

%%%%%%%%%%%%%%%%%%%%%%%% Vectors %%%%%%%%%%%%%%%%%%%%%%%%%%%
\def\aa{{\mathbf{a}}}
\def\bb{{\mathbf{b}}}
\def\cc{{\mathbf{c}}}
\def\dd{{\mathbf{d}}}
\def\ee{{\mathbf{e}}}
\def\ff{{\mathbf{f}}}
\def\gg{{\mathbf{g}}}
\def\hh{{\mathbf{h}}}
\def\nn{{\mathbf{n}}}
\def\rr{{\mathbf{r}}}
\def\ss{{\mathbf{s}}}
\def\tt{{\mathbf{t}}}
\def\uu{{\mathbf{u}}}
\def\vv{{\mathbf{v}}}
\def\ww{{\mathbf{w}}}
\def\xx{{\mathbf{x}}}
\def\yy{{\mathbf{y}}}
\def\zz{{\mathbf{z}}}



\def\aalpha{\boldsymbol{\alpha}}
\def\bbeta{\boldsymbol{\beta}}
\def\ssig{{\boldsymbol{\sigma}}}

%%%%%%%%%%%%%%%%%%%%%%%% Vectors -- Caps %%%%%%%%%%%%%%%%%%%%%%%%%%%
\def\AA{{\mathbf{A}}}
\def\BB{{\mathbf{B}}}
\def\CC{{\mathbf{C}}}
\def\DD{{\mathbf{D}}}
\def\EE{{\mathbf{E}}}
\def\FF{{\mathbf{F}}}
\def\GG{{\mathbf{G}}}
\def\HH{{\mathbf{H}}}
\def\MM{{\mathbf{M}}}
\def\RR{{\mathbf{R}}}
\def\SS{{\mathbf{S}}}
\def\TT{{\mathbf{T}}}
\def\UU{{\mathbf{U}}}
\def\VV{{\mathbf{V}}}
\def\XX{{\mathbf{X}}}
\def\YY{{\mathbf{Y}}}
\def\ZZ{{\mathbf{Z}}}

%%%%%%%%%%%%%%%%%%%%%%%% Caps: Bold %%%%%%%%%%%%%%%%%%%%%%%% 
\newcommand{\B}{\mathbb{B}}
\newcommand{\C}{\mathbb{C}}
\newcommand{\D}{\mathbb{D}}
\newcommand{\E}{\mathbb{E}}
\newcommand{\F}{\mathbb{F}}
\renewcommand{\H}{\mathbb{H}}
% \renewcommand{\K}{\mathbb{K}}
\newcommand{\N}{\mathbb{N}}
\renewcommand{\P}{\mathbb{P}}
\newcommand{\Q}{\mathbb{Q}}
\newcommand{\R}{\mathbb{R}}
\newcommand{\T}{\mathbb{T}}

\newcommand{\TP}{\mathbb{TP}}
\newcommand{\V}{\mathbb{V}}
\newcommand{\X}{\mathbb{X}}
\newcommand{\Z}{\mathbb{Z}}




%%%%%%%%%%%%%%%%%%%%%%%% caps: Calligraphy %%%%%%%%%%%%%%%%%%%%%%%%%%%
\newcommand\cA{{\ensuremath{\mathcal{A}}}\xspace}
\newcommand\cB{{\ensuremath{\mathcal{B}}}\xspace}
\newcommand\cC{{\ensuremath{\mathcal{C}}}\xspace}
\newcommand\cD{{\ensuremath{\mathcal{D}}}\xspace}
\newcommand\cE{{\ensuremath{\mathcal{E}}}\xspace}
\newcommand\cF{{\ensuremath{\mathcal{F}}}\xspace}
\newcommand\cG{{\ensuremath{\mathcal{G}}}\xspace}
\newcommand\cH{{\ensuremath{\mathcal{H}}}\xspace}
\newcommand\cI{{\ensuremath{\mathcal{I}}}\xspace}
\newcommand\cJ{{\ensuremath{\mathcal{J}}}\xspace}
\newcommand\cK{{\ensuremath{\mathcal{K}}}\xspace}
\newcommand\cL{{\ensuremath{\mathcal{L}}}\xspace}
\newcommand\cM{{\ensuremath{\mathcal{M}}}\xspace}
\newcommand\cN{{\ensuremath{\mathcal{N}}}\xspace}
\newcommand\cO{{\ensuremath{\mathcal{O}}}\xspace}
\newcommand\cP{{\ensuremath{\mathcal{P}}}\xspace}
\newcommand\cQ{{\ensuremath{\mathcal{Q}}}\xspace}
\newcommand\cR{{\ensuremath{\mathcal{R}}}\xspace}
\newcommand\cS{{\ensuremath{\mathcal{S}}}\xspace}
\newcommand\cT{{\ensuremath{\mathcal{T}}}\xspace}
\newcommand\cU{{\ensuremath{\mathcal{U}}}\xspace}
\newcommand\cV{{\ensuremath{\mathcal{V}}}\xspace}
\newcommand\cW{{\ensuremath{\mathcal{W}}}\xspace}
\newcommand\cX{{\ensuremath{\mathcal{X}}}\xspace}
\newcommand\cY{{\ensuremath{\mathcal{Y}}}\xspace}
\newcommand\cZ{{\ensuremath{\mathcal{Z}}}\xspace}

%%%%%%%%%%%%%%%%%%%%%%%% Greek %%%%%%%%%%%%%%%%%%%%%%%%%%%
\newcommand{\eps}{\epsilon}
\newcommand{\sig}{\sigma}
\newcommand{\Sig}{\Sigma}

%%%%%%%%%%%%%%%%%%%%%%%% Custom %%%%%%%%%%%%%%%%%%%%%%%%%%%
\DeclareMathOperator{\vertices}{vert}
\DeclareMathOperator{\face}{face}
\DeclareMathOperator{\dist}{dist}
\DeclareMathOperator{\faces}{faces}
\DeclareMathOperator{\acc}{acc}
\DeclareMathOperator{\apker}{apker}
\DeclareMathOperator{\aponb}{APONB}
\DeclareMathOperator{\argmax}{argmax}
\DeclareMathOperator{\argmin}{argmin}
\DeclareMathOperator{\conv}{conv}
\DeclareMathOperator{\diag}{diag}
\DeclareMathOperator{\eval}{eval}
\DeclareMathOperator{\evaluation}{EV}
\DeclareMathOperator{\gr}{g}
\DeclareMathOperator{\gw}{gw}
\DeclareMathOperator{\Ker}{Ker}
\DeclareMathOperator{\lmo}{LMO}
\DeclareMathOperator{\lt}{LT}
\DeclareMathOperator{\ltc}{LTC}
\DeclareMathOperator{\Mat}{Mat}
\DeclareMathOperator{\mse}{MSE}
% \DeclareMathOperator{\rank}{rank}
\DeclareMathOperator{\rmse}{rmse}
\DeclareMathOperator{\mathspan}{span}
\DeclareMathOperator{\spar}{SPAR}
\DeclareMathOperator{\srref}{SRREF}



\newcommand{\zeros}{\ensuremath{\mathbf{0}}}
\newcommand{\ones}{\mathbf{1}}
\newcommand{\oneterm}{\ensuremath{\mathbb{1}}}

\newcommand{\true}{\operatorname{true}}
\newcommand{\false}{\operatorname{false}}

\newcommand{\innp}[1]{\left\langle #1 \right\rangle}
\newcommand{\bdot}[1]{\mathbf{\dot{ #1 }}}
\newcommand{\OPT}{\operatorname{OPT}}
\newcommand{\norm}[1]{\left\| #1 \right\|}
\newcommand{\card}[1]{\left| #1 \right|}
\newcommand{\abs}[1]{\lvert #1 \rvert}

% et. al.
\newcommand{\etal}{\textit{et al}.}

% :=
\newcommand{\defeq}{\stackrel{\mathrm{\scriptscriptstyle def}}{=}}

%%%%%%%%%%%%%%%%%%%%%%%% Lines %%%%%%%%%%%%%%%%%%%%%%%%%%%
\newcommand*{\sepfbox}[1]{%
  \begingroup
    \sbox0{\fbox{#1}}%
    \setlength{\fboxrule}{0pt}%
    \fbox{\unhbox0}%
  \endgroup
}
\newcommand*{\vsepfbox}[1]{%
  \begingroup
    \sbox0{\fbox{#1}}%
    \setlength{\fboxrule}{0pt}%
    \mbox{\kern-\fboxsep\fbox{\unhbox0}\kern-\fboxsep}%
  \endgroup
}
\newcommand*{\vertbar}{\rule[-1ex]{0.5pt}{2.5ex}}
\newcommand*{\horzbar}{\rule[.5ex]{2.5ex}{0.5pt}}

\newenvironment{nscenter}
 {\parskip=0pt\par\nopagebreak\centering}
 {\par\noindent\ignorespacesafterend}

%%%%%%%%%%%%%%%% Theorem setup. %%%%%%%%%%%%%%
\theoremstyle{plain} \numberwithin{equation}{section}
\newtheorem{theorem}{Theorem}[section]
\numberwithin{theorem}{section}
\newtheorem{lemma}[theorem]{Lemma}
\newtheorem{corollary}[theorem]{Corollary}
\newtheorem{proposition}[theorem]{Proposition}
\newtheorem{propoblem}[theorem]{Problem}
\newtheorem{conjecture}[theorem]{Conjecture}
\newtheorem{claim}[theorem]{Claim}
\newtheorem{fact}[theorem]{Fact}
\newtheorem{problem}[theorem]{Problem}
\theoremstyle{definition}
\newtheorem{definition}[theorem]{Definition}
\newtheorem{finalremark}[theorem]{Final Remark}
\newtheorem{remark}[theorem]{Remark}
\newtheorem{example}[theorem]{Example}
\newtheorem{observation}[theorem]{Observation}
\newtheorem{maintheorem}[theorem]{Main Theorem}

\theoremstyle{plain}
\newtheorem{assumption}{Assumption}
\let\origtheassumption\theassumption


%%%%%%%%%%%%%%%%%%%%%%%% Miscelaneous %%%%%%%
\graphicspath{{imgs/}}
\makeatletter
\def\mathcolor#1#{\@mathcolor{#1}}
\def\@mathcolor#1#2#3{%
  \protect\leavevmode
  \begingroup
    \color#1{#2}#3%
  \endgroup
}
\makeatother
\setlength{\parskip}{0.75em}

%%%% Sebastian's macros
\newcommand{\tk}{{\color{red}{\bf tk}}}

%HERE
\iffalse

%For one liner if else statements.
\algnewcommand{\IfThenElse}[3]{% \IfThenElse{<if>}{<then>}{<else>}
  \State \algorithmicif\ #1\ \algorithmicthen\ #2\ \algorithmicelse\ #3}

%Redefine left and right.
\let\originalleft\left
\let\originalright\right
\renewcommand{\left}{\mathopen{}\mathclose\bgroup\originalleft}
\renewcommand{\right}{\aftergroup\egroup\originalright}

\fi

%For computer font.
\newenvironment{myfont}{\fontfamily{<familyname>}\selectfont}{\par}

%Adding the lines in the algorithm between input and output.
\newcommand{\hrulealg}[0]{\vspace{1mm} \hrule \vspace{1mm}}

\usepackage[colorlinks,urlcolor=blue,citecolor=blue,linkcolor=blue]{hyperref}

% 節参照を「第3節」の形式で表示する。
\usepackage{cleveref}
\crefformat{section}{第#2#1#3節}
\Crefformat{section}{第#2#1#3節}

\endinput
